\documentclass[11pt,a4paper]{article}

% ---- Packages ----
\usepackage[utf8]{inputenc}
\usepackage[T1]{fontenc}
\usepackage{amsmath,amssymb,amsthm}
\usepackage{mathtools}
\usepackage[margin=1in,headheight=14pt]{geometry}
\usepackage{hyperref}
\usepackage{enumitem}
\usepackage{xcolor}
\usepackage{tcolorbox}
\usepackage{fancyhdr}
\usepackage{booktabs}

% ---- Hyperref ----
\hypersetup{
    colorlinks=true,
    linkcolor=red,
    citecolor=blue,
    urlcolor=blue
}

% ---- Header ----
\pagestyle{fancy}
\fancyhf{}
\fancyhead[L]{\textit{Expectation Cheatsheet for TD Learning}}
\fancyhead[R]{\thepage}
\renewcommand{\headrulewidth}{0.4pt}

% ---- Theorem Environments ----
\theoremstyle{definition}
\newtheorem{definition}{Definition}[section]

\theoremstyle{plain}
\newtheorem{theorem}{Theorem}[section]
\newtheorem{lemma}[theorem]{Lemma}
\newtheorem{proposition}[theorem]{Proposition}
\newtheorem{corollary}[theorem]{Corollary}

\theoremstyle{remark}
\newtheorem{remark}{Remark}[section]

% ---- Colored Boxes ----
\tcbuselibrary{theorems,skins,breakable}

\newtcolorbox{keyresult}[1][]{
    colback=green!5!white,
    colframe=green!50!black,
    fonttitle=\bfseries,
    title=#1,
    breakable
}

\newtcolorbox{tdconnection}[1][]{
    colback=blue!5!white,
    colframe=blue!50!black,
    fonttitle=\bfseries,
    title=#1,
    breakable
}

% ---- Operators ----
\DeclareMathOperator{\Var}{Var}
\DeclareMathOperator{\Cov}{Cov}
\newcommand{\E}{\mathbb{E}}
\newcommand{\R}{\mathbb{R}}
\newcommand{\Prob}{\mathbb{P}}
\newcommand{\indicator}{\mathbf{1}}

% ---- Title ----
\title{\textbf{Expectation: A Self-Contained Reference\\with Formal Proofs}\\[10pt]
\large Prerequisites for Temporal-Difference Learning}
\author{Companion to \emph{TD Learning: A Comprehensive Mathematical Treatment}}
\date{February 2026}

\begin{document}

\maketitle

\begin{abstract}
This document provides formal, self-contained proofs of every expectation identity and inequality used in the companion article on Temporal-Difference learning. Each result is stated precisely, proved from first principles in the discrete (finite) setting, and linked to its role in TD theory. The reader who masters these results will have the complete probabilistic toolkit needed to follow all proofs in the main article.
\end{abstract}

\tableofcontents
\newpage

%% ============================================================
%% SECTION 1: FOUNDATIONS
%% ============================================================
\section{Foundations}
\label{sec:foundations}

\subsection{Probability Space and Random Variables}
\label{subsec:prob_space}

\begin{definition}[Probability Space]
\label{def:prob_space}
A \emph{probability space} is a triple $(\Omega, \mathcal{F}, \Prob)$ where:
\begin{itemize}[nosep]
    \item $\Omega$ is the sample space (set of all possible outcomes),
    \item $\mathcal{F}$ is a $\sigma$-algebra on $\Omega$ (collection of events closed under complement and countable union),
    \item $\Prob: \mathcal{F} \to [0,1]$ is a probability measure with $\Prob(\Omega) = 1$.
\end{itemize}
\end{definition}

\begin{definition}[Random Variable --- Discrete Case]
\label{def:random_variable}
A \emph{discrete random variable} is a measurable function $X: \Omega \to \R$ that takes values in a countable set $\{x_1, x_2, \ldots\}$. Its \emph{probability mass function} (PMF) is $p_X(x) = \Prob(X = x) = \Prob(\{\omega \in \Omega : X(\omega) = x\})$.
\end{definition}

\subsection{Expectation}
\label{subsec:expectation}

\begin{definition}[Expectation]
\label{def:expectation}
The \emph{expectation} (or \emph{expected value}) of a discrete random variable $X$ with PMF $p_X$ is
\begin{equation}
    \E[X] = \sum_{x} x\, p_X(x),
    \label{eq:expectation_def}
\end{equation}
provided the sum converges absolutely: $\sum_x |x|\, p_X(x) < \infty$.
\end{definition}

\begin{tdconnection}[Connection to TD Learning]
The state-value function $V^\pi(s) = \E_\pi[G_0 \mid S_0 = s]$ is an expectation of the return $G_0 = \sum_{t=0}^\infty \gamma^t R_{t+1}$ over all trajectories from state $s$. Since the MDP has finite state and action spaces and $\gamma < 1$, the return is bounded, so the expectation is always well defined.
\end{tdconnection}

%% ============================================================
%% SECTION 2: LINEARITY
%% ============================================================
\section{Linearity of Expectation}
\label{sec:linearity}

\begin{keyresult}[Linearity of Expectation]
\begin{theorem}[Linearity]
\label{thm:linearity}
For any discrete random variables $X$ and $Y$ with finite expectations, and constants $a, b \in \R$:
\begin{equation}
    \E[aX + bY] = a\,\E[X] + b\,\E[Y].
    \label{eq:linearity}
\end{equation}
This holds regardless of whether $X$ and $Y$ are independent.
\end{theorem}
\end{keyresult}

\begin{proof}
Let $X$ take values $\{x_i\}$ and $Y$ take values $\{y_j\}$. Denote the joint PMF by $p_{X,Y}(x_i, y_j) = \Prob(X = x_i, Y = y_j)$. Then:
\begin{align}
    \E[aX + bY]
    &= \sum_{i}\sum_{j} (a x_i + b y_j)\, p_{X,Y}(x_i, y_j) \nonumber\\
    &= a \sum_{i}\sum_{j} x_i\, p_{X,Y}(x_i, y_j) + b \sum_{i}\sum_{j} y_j\, p_{X,Y}(x_i, y_j) \nonumber\\
    &= a \sum_{i} x_i \underbrace{\sum_{j} p_{X,Y}(x_i, y_j)}_{= p_X(x_i)} + b \sum_{j} y_j \underbrace{\sum_{i} p_{X,Y}(x_i, y_j)}_{= p_Y(y_j)} \nonumber\\
    &= a \sum_{i} x_i\, p_X(x_i) + b \sum_{j} y_j\, p_Y(y_j) \nonumber\\
    &= a\,\E[X] + b\,\E[Y]. \label{eq:linearity_proof}
\end{align}
The marginal sums $\sum_j p_{X,Y}(x_i, y_j) = p_X(x_i)$ and $\sum_i p_{X,Y}(x_i, y_j) = p_Y(y_j)$ follow from the definition of marginal distributions. No independence assumption is needed.
\end{proof}

\begin{corollary}[Finite Sums]
\label{cor:linearity_sum}
For any random variables $X_1, \ldots, X_n$ with finite expectations and constants $a_1, \ldots, a_n$:
\begin{equation}
    \E\!\left[\sum_{k=1}^n a_k X_k\right] = \sum_{k=1}^n a_k\,\E[X_k].
    \label{eq:linearity_sum}
\end{equation}
\end{corollary}

\begin{proof}
Apply Theorem~\ref{thm:linearity} inductively $n-1$ times.
\end{proof}

\begin{tdconnection}[Connection to TD Learning]
Linearity is used to split the return $G_t = R_{t+1} + \gamma G_{t+1}$ inside expectations:
\[
    \E_\pi[G_t \mid S_t = s] = \E_\pi[R_{t+1} \mid S_t = s] + \gamma\,\E_\pi[G_{t+1} \mid S_t = s].
\]
This is the first step of every Bellman equation derivation.
\end{tdconnection}

%% ============================================================
%% SECTION 3: CONDITIONAL EXPECTATION
%% ============================================================
\section{Conditional Expectation}
\label{sec:conditional}

\subsection{Definition and Basic Properties}
\label{subsec:cond_def}

\begin{definition}[Conditional Expectation Given an Event]
\label{def:cond_exp_event}
For a discrete random variable $X$ and an event $B$ with $\Prob(B) > 0$:
\begin{equation}
    \E[X \mid B] = \sum_x x\, \Prob(X = x \mid B) = \sum_x x\, \frac{\Prob(X = x,\, B)}{\Prob(B)}.
    \label{eq:cond_exp_event}
\end{equation}
\end{definition}

\begin{definition}[Conditional Expectation Given a Random Variable]
\label{def:cond_exp_rv}
For discrete random variables $X$ and $Y$, the \emph{conditional expectation of $X$ given $Y$} is the random variable $\E[X \mid Y]$ defined by
\begin{equation}
    \E[X \mid Y](\omega) = \E[X \mid Y = Y(\omega)] = \sum_x x\, \Prob(X = x \mid Y = Y(\omega)).
    \label{eq:cond_exp_rv}
\end{equation}
That is, $\E[X \mid Y]$ is a function of $Y$: it assigns to each outcome $\omega$ the conditional mean of $X$ given the observed value of $Y$.
\end{definition}

\begin{tdconnection}[Connection to TD Learning]
The value function $V^\pi(s) = \E_\pi[G_t \mid S_t = s]$ is a conditional expectation given a specific state. The notation $\E_\pi[\delta_t^\pi \mid S_t = s]$ conditions on the current state $s$ and averages over the randomness of $A_t$, $R_{t+1}$, and $S_{t+1}$.
\end{tdconnection}

\subsection{Linearity of Conditional Expectation}
\label{subsec:cond_linearity}

\begin{theorem}[Linearity of Conditional Expectation]
\label{thm:cond_linearity}
For discrete random variables $X$, $Y$, $Z$ and constants $a, b \in \R$:
\begin{equation}
    \E[aX + bY \mid Z] = a\,\E[X \mid Z] + b\,\E[Y \mid Z].
    \label{eq:cond_linearity}
\end{equation}
\end{theorem}

\begin{proof}
Fix $Z = z$ with $\Prob(Z = z) > 0$.

\textbf{Step 1: Expand the definition of conditional expectation.}
The conditional expectation of $aX + bY$ given $Z = z$ is computed by summing over all joint values $(x, y)$:
\begin{equation}
    \E[aX + bY \mid Z = z]
    = \sum_{x,y} (ax + by)\, \Prob(X = x, Y = y \mid Z = z).
    \label{eq:cond_lin_step1}
\end{equation}

\textbf{Step 2: Distribute the sum.}
Since the sum $\sum_{x,y}$ is a finite sum and $(ax + by) = ax + by$, we can split it into two separate sums:
\begin{equation}
    = a \sum_{x,y} x\, \Prob(X = x, Y = y \mid Z = z)
    \;+\; b \sum_{x,y} y\, \Prob(X = x, Y = y \mid Z = z).
    \label{eq:cond_lin_step2}
\end{equation}

\textbf{Step 3: Marginalize --- reduce $\sum_{x,y}$ to $\sum_x$ (and $\sum_y$).}
This is the key step. Consider the first sum. Since $x$ does \emph{not} depend on $y$, we can rearrange the double sum by summing over $y$ first, with $x$ held fixed:
\begin{align}
    \sum_{x,y} x\, \Prob(X = x, Y = y \mid Z = z)
    &= \sum_x x \underbrace{\sum_y \Prob(X = x, Y = y \mid Z = z)}_{\text{sum over all possible } y}.
    \label{eq:cond_lin_step3a}
\end{align}
The inner sum over $y$ is the \textbf{marginalization} (law of total probability applied conditionally). Since the events $\{Y = y\}$ for all $y$ form a partition of the sample space:
\begin{align}
    \sum_y \Prob(X = x, Y = y \mid Z = z)
    &= \Prob\!\left(\bigcup_y \{X = x, Y = y\} \;\middle|\; Z = z\right) \nonumber\\
    &= \Prob(X = x \mid Z = z).
    \label{eq:cond_lin_marginal}
\end{align}
The union $\bigcup_y \{X = x, Y = y\}$ covers every possible value of $Y$ while keeping $X = x$ fixed, so it equals the event $\{X = x\}$.

Substituting~\eqref{eq:cond_lin_marginal} back into~\eqref{eq:cond_lin_step3a}:
\begin{equation}
    \sum_{x,y} x\, \Prob(X = x, Y = y \mid Z = z) = \sum_x x\, \Prob(X = x \mid Z = z).
    \label{eq:cond_lin_step3b}
\end{equation}

By the \emph{same reasoning} applied to the second sum (sum over $x$ first, marginalize out $X$):
\begin{equation}
    \sum_{x,y} y\, \Prob(X = x, Y = y \mid Z = z) = \sum_y y\, \Prob(Y = y \mid Z = z).
    \label{eq:cond_lin_step3c}
\end{equation}

\textbf{Step 4: Recognize the conditional expectations.}
Substituting~\eqref{eq:cond_lin_step3b} and~\eqref{eq:cond_lin_step3c} into~\eqref{eq:cond_lin_step2}:
\begin{align}
    \E[aX + bY \mid Z = z]
    &= a \sum_x x\, \Prob(X = x \mid Z = z) + b \sum_y y\, \Prob(Y = y \mid Z = z) \nonumber\\
    &= a\,\E[X \mid Z = z] + b\,\E[Y \mid Z = z].
    \label{eq:cond_lin_step4}
\end{align}

Since this holds for every value $z$ with $\Prob(Z = z) > 0$, we conclude as an identity of random variables:
\begin{equation}
    \E[aX + bY \mid Z] = a\,\E[X \mid Z] + b\,\E[Y \mid Z]. \qedhere
\end{equation}
\end{proof}

\subsection{Pulling Out Known Factors}
\label{subsec:pull_out}

\begin{keyresult}[Pull-Out Property]
\begin{theorem}[Pull-Out Property]
\label{thm:pull_out}
If $h(Z)$ is a function of $Z$ (i.e., $h(Z)$ is $\sigma(Z)$-measurable), then:
\begin{equation}
    \E[h(Z) \cdot X \mid Z] = h(Z) \cdot \E[X \mid Z].
    \label{eq:pull_out}
\end{equation}
\end{theorem}
\end{keyresult}

\begin{proof}
Fix $Z = z$. Then $h(Z) = h(z)$ is a constant under this conditioning:
\begin{align}
    \E[h(Z) \cdot X \mid Z = z]
    &= \sum_x h(z) \cdot x \cdot \Prob(X = x \mid Z = z) \nonumber\\
    &= h(z) \sum_x x\, \Prob(X = x \mid Z = z) \nonumber\\
    &= h(z) \cdot \E[X \mid Z = z]. \label{eq:pull_out_proof}
\end{align}
Since this holds for every $z$, the result follows as an identity of random variables.
\end{proof}

\begin{tdconnection}[Connection to TD Learning]
When proving that the TD error has zero conditional mean, we condition on $S_t = s$ and pull out $V(S_t) = V(s)$:
\[
    \E[\delta_t \mid S_t = s] = \E[R_{t+1} + \gamma V(S_{t+1}) \mid S_t = s] - V(s).
\]
Here $V(s)$ exits the expectation because it is fully determined by $S_t$.
\end{tdconnection}

%% ============================================================
%% SECTION 4: EXPANDING OVER ACTIONS AND TRANSITIONS
%% ============================================================
\section{Expanding Expectations in an MDP}
\label{sec:expanding}

This section formalizes the mechanical step that appears in virtually every derivation in the TD article.

\begin{theorem}[MDP Expectation Expansion]
\label{thm:mdp_expansion}
Let $(\mathcal{S}, \mathcal{A}, P, r, \gamma)$ be a finite MDP with policy $\pi$. For any function $f: \mathcal{A} \times \mathcal{S} \times \R \to \R$:
\begin{equation}
    \E_\pi[f(A_t, S_{t+1}, R_{t+1}) \mid S_t = s] = \sum_{a \in \mathcal{A}} \pi(a|s) \sum_{s' \in \mathcal{S}} P(s'|s,a)\, f(a, s', r(s,a,s')).
    \label{eq:mdp_expansion}
\end{equation}
\end{theorem}

\begin{proof}
By the definition of conditional expectation and the MDP dynamics:
\begin{align}
    &\E_\pi[f(A_t, S_{t+1}, R_{t+1}) \mid S_t = s] \nonumber\\
    &\quad= \sum_{a \in \mathcal{A}} \sum_{s' \in \mathcal{S}} f(a, s', r(s,a,s'))\, \Prob(A_t = a, S_{t+1} = s' \mid S_t = s). \label{eq:expand_step1}
\end{align}
By the chain rule for conditional probability:
\begin{align}
    \Prob(A_t = a, S_{t+1} = s' \mid S_t = s)
    &= \Prob(A_t = a \mid S_t = s) \cdot \Prob(S_{t+1} = s' \mid S_t = s, A_t = a) \nonumber\\
    &= \pi(a|s) \cdot P(s'|s,a). \label{eq:expand_step2}
\end{align}
Substituting~\eqref{eq:expand_step2} into~\eqref{eq:expand_step1} yields~\eqref{eq:mdp_expansion}.
\end{proof}

\begin{corollary}[Bellman Expectation --- Expectation Form]
\label{cor:bellman_expansion}
Setting $f(a, s', r) = r + \gamma V(s')$ in Theorem~\ref{thm:mdp_expansion}:
\begin{equation}
    \E_\pi[R_{t+1} + \gamma V(S_{t+1}) \mid S_t = s] = \sum_a \pi(a|s) \sum_{s'} P(s'|s,a)\bigl[r(s,a,s') + \gamma V(s')\bigr] = (\mathcal{T}^\pi V)(s).
    \label{eq:bellman_expansion}
\end{equation}
\end{corollary}

\begin{proof}
Direct substitution into~\eqref{eq:mdp_expansion}.
\end{proof}

%% ============================================================
%% SECTION 5: LAW OF TOTAL EXPECTATION
%% ============================================================
\section{Law of Total Expectation (Tower Property)}
\label{sec:tower}

\begin{keyresult}[Law of Total Expectation]
\begin{theorem}[Law of Total Expectation]
\label{thm:tower}
For discrete random variables $X$ and $Y$ with $\E[|X|] < \infty$:
\begin{equation}
    \E[X] = \E\!\left[\E[X \mid Y]\right] = \sum_y \E[X \mid Y = y]\, \Prob(Y = y).
    \label{eq:tower}
\end{equation}
\end{theorem}
\end{keyresult}

\begin{proof}
Expand the right-hand side:
\begin{align}
    \sum_y \E[X \mid Y = y]\, \Prob(Y = y)
    &= \sum_y \left(\sum_x x\, \Prob(X = x \mid Y = y)\right) \Prob(Y = y) \nonumber\\
    &= \sum_y \sum_x x\, \Prob(X = x \mid Y = y)\, \Prob(Y = y) \nonumber\\
    &= \sum_y \sum_x x\, \frac{\Prob(X = x, Y = y)}{\Prob(Y = y)} \cdot \Prob(Y = y) \nonumber\\
    &= \sum_y \sum_x x\, \Prob(X = x, Y = y) \nonumber\\
    &= \sum_x x \sum_y \Prob(X = x, Y = y) \nonumber\\
    &= \sum_x x\, \Prob(X = x) = \E[X]. \label{eq:tower_proof}
\end{align}
The interchange of summation is justified by absolute convergence.
\end{proof}

\begin{theorem}[Generalized Tower Property]
\label{thm:tower_general}
If $\mathcal{G} \subseteq \mathcal{H}$ are two $\sigma$-algebras (i.e., $\mathcal{G}$ contains ``less information'' than $\mathcal{H}$), then:
\begin{equation}
    \E\!\left[\E[X \mid \mathcal{H}] \;\middle|\; \mathcal{G}\right] = \E[X \mid \mathcal{G}].
    \label{eq:tower_general}
\end{equation}
In particular, for random variables with $\sigma(Z) \subseteq \sigma(Y, Z)$:
\begin{equation}
    \E\!\left[\E[X \mid Y, Z] \;\middle|\; Z\right] = \E[X \mid Z].
    \label{eq:tower_yz}
\end{equation}
\end{theorem}

\begin{proof}
We prove~\eqref{eq:tower_yz} in the discrete case. Fix $Z = z$:
\begin{align}
    \E\!\left[\E[X \mid Y, Z] \;\middle|\; Z = z\right]
    &= \sum_y \E[X \mid Y = y, Z = z]\, \Prob(Y = y \mid Z = z) \nonumber\\
    &= \sum_y \left(\sum_x x\, \Prob(X = x \mid Y = y, Z = z)\right) \Prob(Y = y \mid Z = z) \nonumber\\
    &= \sum_x x \sum_y \Prob(X = x \mid Y = y, Z = z)\, \Prob(Y = y \mid Z = z) \nonumber\\
    &= \sum_x x \sum_y \frac{\Prob(X = x, Y = y, Z = z)}{\Prob(Y = y, Z = z)} \cdot \frac{\Prob(Y = y, Z = z)}{\Prob(Z = z)} \nonumber\\
    &= \sum_x x \sum_y \frac{\Prob(X = x, Y = y, Z = z)}{\Prob(Z = z)} \nonumber\\
    &= \sum_x x\, \frac{\Prob(X = x, Z = z)}{\Prob(Z = z)} \nonumber\\
    &= \sum_x x\, \Prob(X = x \mid Z = z) = \E[X \mid Z = z]. \label{eq:tower_general_proof}
\end{align}
\end{proof}

\begin{tdconnection}[Connection to TD Learning]
The Bellman equation derivation uses the tower property with the Markov property. Starting from $V^\pi(s) = \E_\pi[R_1 + \gamma G_1 \mid S_0 = s]$:
\[
    \E_\pi[G_1 \mid S_0 = s] = \E_\pi\!\left[\E_\pi[G_1 \mid S_1] \;\middle|\; S_0 = s\right] = \E_\pi[V^\pi(S_1) \mid S_0 = s].
\]
The inner conditioning on $S_1$ uses the Markov property ($S_1$ screens off the future from $S_0$), and the outer averaging over $S_1$ given $S_0$ uses the tower property.
\end{tdconnection}

%% ============================================================
%% SECTION 6: VARIANCE
%% ============================================================
\section{Variance and Conditional Variance}
\label{sec:variance}

\subsection{Definitions}
\label{subsec:var_def}

\begin{definition}[Variance]
\label{def:variance}
The \emph{variance} of a random variable $X$ with $\E[X^2] < \infty$ is
\begin{equation}
    \Var(X) = \E\!\left[(X - \E[X])^2\right].
    \label{eq:var_def}
\end{equation}
\end{definition}

\begin{theorem}[Computational Formula for Variance]
\label{thm:var_formula}
\begin{equation}
    \Var(X) = \E[X^2] - (\E[X])^2.
    \label{eq:var_formula}
\end{equation}
\end{theorem}

\begin{proof}
Let $\mu = \E[X]$. Expanding the square:
\begin{align}
    \Var(X) &= \E[(X - \mu)^2] \nonumber\\
    &= \E[X^2 - 2\mu X + \mu^2] \nonumber\\
    &= \E[X^2] - 2\mu\,\E[X] + \mu^2 \tag{linearity} \\
    &= \E[X^2] - 2\mu^2 + \mu^2 \nonumber\\
    &= \E[X^2] - \mu^2 = \E[X^2] - (\E[X])^2. \label{eq:var_formula_proof}
\end{align}
\end{proof}

\begin{definition}[Conditional Variance]
\label{def:cond_var}
The \emph{conditional variance} of $X$ given $Y$ is the random variable
\begin{equation}
    \Var(X \mid Y) = \E\!\left[(X - \E[X \mid Y])^2 \;\middle|\; Y\right] = \E[X^2 \mid Y] - (\E[X \mid Y])^2.
    \label{eq:cond_var}
\end{equation}
\end{definition}

\begin{proof}[Proof of the computational formula~\eqref{eq:cond_var}]
Fix $Y = y$ and let $\mu_y = \E[X \mid Y = y]$. Applying the same expansion as in Theorem~\ref{thm:var_formula} to the conditional expectation:
\begin{align}
    \Var(X \mid Y = y) &= \E[(X - \mu_y)^2 \mid Y = y] \nonumber\\
    &= \E[X^2 \mid Y = y] - 2\mu_y\,\E[X \mid Y = y] + \mu_y^2 \nonumber\\
    &= \E[X^2 \mid Y = y] - \mu_y^2. \label{eq:cond_var_proof}
\end{align}
\end{proof}

\subsection{Law of Total Variance}
\label{subsec:total_var}

\begin{keyresult}[Law of Total Variance]
\begin{theorem}[Law of Total Variance]
\label{thm:total_variance}
For random variables $X$ and $Y$ with $\E[X^2] < \infty$:
\begin{equation}
    \Var(X) = \E\!\left[\Var(X \mid Y)\right] + \Var\!\left(\E[X \mid Y]\right).
    \label{eq:total_variance}
\end{equation}
\end{theorem}
\end{keyresult}

\begin{proof}
We compute each term on the right-hand side.

\textbf{First term.} Using~\eqref{eq:cond_var}:
\begin{align}
    \E[\Var(X \mid Y)]
    &= \E\!\left[\E[X^2 \mid Y] - (\E[X \mid Y])^2\right] \nonumber\\
    &= \E[\E[X^2 \mid Y]] - \E[(\E[X \mid Y])^2] \nonumber\\
    &= \E[X^2] - \E[(\E[X \mid Y])^2], \label{eq:total_var_first}
\end{align}
where the last equality uses the tower property $\E[\E[X^2 \mid Y]] = \E[X^2]$.

\textbf{Second term.} Let $\mu = \E[X]$. By the computational formula for variance:
\begin{align}
    \Var(\E[X \mid Y])
    &= \E[(\E[X \mid Y])^2] - (\E[\E[X \mid Y]])^2 \nonumber\\
    &= \E[(\E[X \mid Y])^2] - (\E[X])^2 \nonumber\\
    &= \E[(\E[X \mid Y])^2] - \mu^2. \label{eq:total_var_second}
\end{align}

\textbf{Sum.} Adding~\eqref{eq:total_var_first} and~\eqref{eq:total_var_second}:
\begin{align}
    \E[\Var(X \mid Y)] + \Var(\E[X \mid Y])
    &= \bigl(\E[X^2] - \E[(\E[X \mid Y])^2]\bigr) + \bigl(\E[(\E[X \mid Y])^2] - \mu^2\bigr) \nonumber\\
    &= \E[X^2] - \mu^2 = \Var(X). \label{eq:total_var_sum}
\end{align}
\end{proof}

\begin{tdconnection}[Connection to TD Learning]
This is the central result for the bias-variance analysis. The article applies it with $X = G_t$ and $Y = (A_t, S_{t+1})$:
\[
    \Var_\pi(G_t \mid S_t) = \underbrace{\Var_\pi(R_{t+1} + \gamma V^\pi(S_{t+1}) \mid S_t)}_{\text{one-step variance}} + \gamma^2 \underbrace{\E_\pi[\Var_\pi(G_{t+1} \mid S_{t+1}) \mid S_t]}_{\text{future variance} \;\geq\; 0}.
\]
The TD(0) target $R_{t+1} + \gamma V(S_{t+1})$ captures only the first term, eliminating the non-negative future variance entirely. This is why TD has lower variance than Monte Carlo.
\end{tdconnection}

%% ============================================================
%% SECTION 7: FILTRATIONS AND CONDITIONAL EXPECTATION
%% ============================================================
\section{Filtrations}
\label{sec:filtrations}

\begin{definition}[Filtration]
\label{def:filtration}
A \emph{filtration} on $(\Omega, \mathcal{F}, \Prob)$ is an increasing sequence of sub-$\sigma$-algebras:
\begin{equation}
    \mathcal{F}_0 \subseteq \mathcal{F}_1 \subseteq \mathcal{F}_2 \subseteq \cdots \subseteq \mathcal{F}.
    \label{eq:filtration}
\end{equation}
\end{definition}

\begin{definition}[Natural Filtration of an MDP]
\label{def:mdp_filtration}
The \emph{natural filtration} of the MDP trajectory is
\begin{equation}
    \mathcal{F}_t = \sigma(S_0, A_0, R_1, S_1, A_1, R_2, \ldots, S_t),
    \label{eq:mdp_filtration}
\end{equation}
the $\sigma$-algebra generated by all observations up to and including time $t$.
\end{definition}

\begin{remark}[Interpretation]
$\mathcal{F}_t$ represents ``all information available at time $t$.'' Saying that a random variable $Z$ is $\mathcal{F}_t$-measurable means $Z$ can be computed from $(S_0, A_0, R_1, \ldots, S_t)$ alone. In particular:
\begin{itemize}[nosep]
    \item $V(S_t)$ is $\mathcal{F}_t$-measurable (it depends only on $S_t$, which is in $\mathcal{F}_t$).
    \item $R_{t+1}$ and $S_{t+1}$ are \emph{not} $\mathcal{F}_t$-measurable (they are ``future'' information).
\end{itemize}
\end{remark}

\begin{proposition}[Conditioning on $\mathcal{F}_t$ Refines Conditioning on $S_t$]
\label{prop:filtration_refine}
In an MDP with the Markov property, for any function of future observations $f(S_{t+1}, R_{t+1}, S_{t+2}, \ldots)$:
\begin{equation}
    \E[f(S_{t+1}, R_{t+1}, \ldots) \mid \mathcal{F}_t] = \E[f(S_{t+1}, R_{t+1}, \ldots) \mid S_t].
    \label{eq:markov_filtration}
\end{equation}
\end{proposition}

\begin{proof}
The Markov property states that the conditional distribution of the future trajectory given the entire history $\mathcal{F}_t$ depends only on $S_t$:
\begin{equation}
    \Prob(S_{t+1} = s', R_{t+1} = r, \ldots \mid \mathcal{F}_t) = \Prob(S_{t+1} = s', R_{t+1} = r, \ldots \mid S_t).
\end{equation}
Since conditional expectations are determined by conditional distributions,~\eqref{eq:markov_filtration} follows immediately.
\end{proof}

\begin{tdconnection}[Connection to TD Learning]
This is why the TD convergence proofs can freely switch between conditioning on the filtration $\mathcal{F}_t$ (needed for stochastic approximation theory) and conditioning on $S_t$ (more convenient for computation). The Markov property guarantees these are equivalent.
\end{tdconnection}

%% ============================================================
%% SECTION 8: MARTINGALE DIFFERENCE SEQUENCES
%% ============================================================
\section{Martingale Difference Sequences}
\label{sec:mds}

\subsection{Definitions}
\label{subsec:mds_def}

\begin{definition}[Martingale]
\label{def:martingale}
A sequence of random variables $\{M_n\}_{n \geq 0}$ is a \emph{martingale} with respect to a filtration $\{\mathcal{F}_n\}$ if for all $n \geq 0$:
\begin{enumerate}[nosep]
    \item $M_n$ is $\mathcal{F}_n$-measurable (adapted),
    \item $\E[|M_n|] < \infty$ (integrable),
    \item $\E[M_{n+1} \mid \mathcal{F}_n] = M_n$ (the ``fair game'' property).
\end{enumerate}
\end{definition}

\begin{definition}[Martingale Difference Sequence]
\label{def:mds}
A sequence $\{w_n\}_{n \geq 1}$ is a \emph{martingale difference sequence} (MDS) with respect to $\{\mathcal{F}_n\}$ if for all $n \geq 1$:
\begin{enumerate}[nosep]
    \item $w_n$ is $\mathcal{F}_n$-measurable,
    \item $\E[|w_n|] < \infty$,
    \item $\E[w_n \mid \mathcal{F}_{n-1}] = 0$.
\end{enumerate}
\end{definition}

\subsection{Fundamental Properties}
\label{subsec:mds_props}

\begin{theorem}[MDS Implies Zero Unconditional Mean]
\label{thm:mds_zero_mean}
If $\{w_n\}$ is an MDS, then $\E[w_n] = 0$ for all $n$.
\end{theorem}

\begin{proof}
By the tower property (Theorem~\ref{thm:tower}):
\begin{equation}
    \E[w_n] = \E[\E[w_n \mid \mathcal{F}_{n-1}]] = \E[0] = 0. \label{eq:mds_zero_proof}
\end{equation}
\end{proof}

\begin{theorem}[MDS Partial Sums Form a Martingale]
\label{thm:mds_martingale}
If $\{w_n\}_{n \geq 1}$ is an MDS with respect to $\{\mathcal{F}_n\}$, then $M_n = \sum_{k=1}^n w_k$ is a martingale with respect to $\{\mathcal{F}_n\}$ (with $M_0 = 0$).
\end{theorem}

\begin{proof}
We verify the three martingale conditions.

\textit{Adapted:} Each $w_k$ is $\mathcal{F}_k$-measurable, hence $\mathcal{F}_n$-measurable for $k \leq n$. So $M_n = \sum_{k=1}^n w_k$ is $\mathcal{F}_n$-measurable.

\textit{Integrable:} $\E[|M_n|] \leq \sum_{k=1}^n \E[|w_k|] < \infty$ by the triangle inequality.

\textit{Fair game:}
\begin{align}
    \E[M_{n+1} \mid \mathcal{F}_n]
    &= \E\!\left[\sum_{k=1}^{n+1} w_k \;\middle|\; \mathcal{F}_n\right] \nonumber\\
    &= \sum_{k=1}^n \E[w_k \mid \mathcal{F}_n] + \E[w_{n+1} \mid \mathcal{F}_n] \nonumber\\
    &= \sum_{k=1}^n w_k + 0 = M_n. \label{eq:mds_martingale_proof}
\end{align}
Here we used: (a) linearity of conditional expectation; (b) $w_k$ is $\mathcal{F}_n$-measurable for $k \leq n$, so $\E[w_k \mid \mathcal{F}_n] = w_k$; (c) the MDS condition $\E[w_{n+1} \mid \mathcal{F}_n] = 0$.
\end{proof}

\begin{theorem}[Uncorrelated Increments of an MDS]
\label{thm:mds_uncorrelated}
If $\{w_n\}$ is an MDS with $\E[w_n^2] < \infty$ for all $n$, then for $m \neq n$:
\begin{equation}
    \E[w_m w_n] = 0.
    \label{eq:mds_uncorrelated}
\end{equation}
Consequently, $\Var(M_n) = \sum_{k=1}^n \E[w_k^2]$.
\end{theorem}

\begin{proof}
Without loss of generality, assume $m < n$. Then $w_m$ is $\mathcal{F}_{n-1}$-measurable. By the pull-out property (Theorem~\ref{thm:pull_out}) and the tower property:
\begin{align}
    \E[w_m w_n]
    &= \E[\E[w_m w_n \mid \mathcal{F}_{n-1}]] \nonumber\\
    &= \E[w_m \cdot \E[w_n \mid \mathcal{F}_{n-1}]] \nonumber\\
    &= \E[w_m \cdot 0] = 0. \label{eq:mds_uncorrelated_proof}
\end{align}
For the variance of the partial sum:
\begin{align}
    \Var(M_n) &= \E[M_n^2] - (\E[M_n])^2 = \E[M_n^2] \qquad \text{(since } \E[M_n] = 0\text{)} \nonumber\\
    &= \E\!\left[\left(\sum_{k=1}^n w_k\right)^2\right]
    = \sum_{k=1}^n \E[w_k^2] + 2\sum_{j < k} \underbrace{\E[w_j w_k]}_{= 0}
    = \sum_{k=1}^n \E[w_k^2]. \label{eq:mds_var_proof}
\end{align}
\end{proof}

\begin{tdconnection}[Connection to TD Learning]
The TD error under the true value function is an MDS:
\[
    \E_\pi[\delta_t^\pi \mid \mathcal{F}_t] = \E_\pi[R_{t+1} + \gamma V^\pi(S_{t+1}) - V^\pi(S_t) \mid S_t] = V^\pi(S_t) - V^\pi(S_t) = 0.
\]
This is \emph{the} central fact enabling all convergence proofs. In the stochastic approximation formulation, the noise term $w_k(s) = R_{n_k+1} + \gamma V_k(S_{n_k+1}) - (\mathcal{T}^\pi V_k)(s)$ satisfies $\E[w_k(s) \mid \mathcal{F}_k] = 0$. The uncorrelated increments property (Theorem~\ref{thm:mds_uncorrelated}) ensures accumulated noise doesn't grow too fast, which is essential for the Robbins-Monro conditions.
\end{tdconnection}

%% ============================================================
%% SECTION 9: JENSEN'S INEQUALITY
%% ============================================================
\section{Jensen's Inequality}
\label{sec:jensen}

\begin{definition}[Convex Function]
\label{def:convex}
A function $\varphi: \R \to \R$ is \emph{convex} if for all $x, y \in \R$ and $\lambda \in [0,1]$:
\begin{equation}
    \varphi(\lambda x + (1-\lambda) y) \leq \lambda\, \varphi(x) + (1-\lambda)\, \varphi(y).
    \label{eq:convex_def}
\end{equation}
\end{definition}

\begin{lemma}[Supporting Hyperplane / Tangent Line]
\label{lem:supporting}
If $\varphi$ is convex, then for any $c \in \R$, there exists a constant $m \in \R$ (a subgradient) such that:
\begin{equation}
    \varphi(x) \geq \varphi(c) + m(x - c) \quad \text{for all } x \in \R.
    \label{eq:supporting}
\end{equation}
\end{lemma}

\begin{proof}
Define $m = \sup_{y < c} \frac{\varphi(c) - \varphi(y)}{c - y}$. By convexity, the slopes of secant lines through $(c, \varphi(c))$ are non-decreasing. For any $x > c$, the secant slope from $c$ to $x$ is at least $m$:
\begin{equation}
    \frac{\varphi(x) - \varphi(c)}{x - c} \geq m,
\end{equation}
which gives $\varphi(x) \geq \varphi(c) + m(x - c)$. For $x < c$, the secant slope from $x$ to $c$ is at most $m$:
\begin{equation}
    \frac{\varphi(c) - \varphi(x)}{c - x} \leq m,
\end{equation}
which rearranges to $\varphi(x) \geq \varphi(c) + m(x - c)$. The case $x = c$ is trivial.
\end{proof}

\begin{keyresult}[Jensen's Inequality]
\begin{theorem}[Jensen's Inequality]
\label{thm:jensen}
If $\varphi$ is convex and $\E[|X|] < \infty$, then:
\begin{equation}
    \varphi\!\left(\E[X]\right) \leq \E[\varphi(X)].
    \label{eq:jensen}
\end{equation}
Conditional version: $\varphi\!\left(\E[X \mid Y]\right) \leq \E[\varphi(X) \mid Y]$.
\end{theorem}
\end{keyresult}

\begin{proof}
Let $c = \E[X]$. By Lemma~\ref{lem:supporting}, there exists $m$ such that $\varphi(X) \geq \varphi(c) + m(X - c)$ pointwise. Taking expectations:
\begin{align}
    \E[\varphi(X)]
    &\geq \E[\varphi(c) + m(X - c)] \nonumber\\
    &= \varphi(c) + m\bigl(\E[X] - c\bigr) \tag{linearity}\\
    &= \varphi(c) + m \cdot 0 = \varphi(c) = \varphi(\E[X]). \label{eq:jensen_proof}
\end{align}

For the conditional version, fix $Y = y$ and apply the unconditional Jensen's inequality to the conditional distribution of $X$ given $Y = y$. Since this holds for every $y$, it holds as a random variable inequality.
\end{proof}

\begin{tdconnection}[Connection to TD Learning]
Jensen's inequality is used in Appendix B of the article to prove the contraction property $\|\mathbf{P}_\pi \mathbf{v}\|_{d^\pi} \leq \|\mathbf{v}\|_{d^\pi}$. The key step: $(\mathbf{P}_\pi \mathbf{v})(s') = \sum_s \frac{d^\pi(s) P_\pi(s'|s)}{d^\pi(s')} v(s)$ is a weighted average of $v(s)$ with weights summing to $1$ (a conditional expectation). Since $\varphi(x) = x^2$ is convex:
\[
    \left(\sum_s w_s\, v(s)\right)^2 \leq \sum_s w_s\, v(s)^2.
\]
This proves $\mathbf{P}_\pi$ is non-expansive in the $d^\pi$-weighted norm, which in turn shows the matrix $\mathbf{A}$ in the linear TD fixed-point equation is negative definite.
\end{tdconnection}

%% ============================================================
%% SECTION 10: NOISE DECOMPOSITION PATTERN
%% ============================================================
\section{The Stochastic Approximation Noise Decomposition}
\label{sec:noise}

This section collects the key expectation facts into the single pattern that drives all TD convergence proofs.

\begin{theorem}[TD Noise Decomposition]
\label{thm:noise_decomp}
In a finite MDP under policy $\pi$, define the noise at state visit $k$ of state $s$:
\begin{equation}
    w_k(s) = R_{n_k+1} + \gamma V_k(S_{n_k+1}) - (\mathcal{T}^\pi V_k)(s),
    \label{eq:noise_def}
\end{equation}
where $n_k$ is the time of the $k$-th visit to $s$. Then:
\begin{enumerate}[nosep]
    \item $\E[w_k(s) \mid \mathcal{F}_k] = 0$ \hfill (martingale difference)
    \item $\E[w_k(s)^2 \mid \mathcal{F}_k] \leq C(1 + \|V_k\|_\infty^2)$ for some constant $C$ \hfill (bounded variance)
\end{enumerate}
\end{theorem}

\begin{proof}
\textbf{Part 1.} Conditioning on $\mathcal{F}_k$ (which includes $V_k$ and $S_{n_k} = s$), and using Proposition~\ref{prop:filtration_refine} to reduce conditioning on $\mathcal{F}_k$ to conditioning on $S_{n_k} = s$:
\begin{align}
    \E[w_k(s) \mid \mathcal{F}_k]
    &= \E[R_{n_k+1} + \gamma V_k(S_{n_k+1}) \mid S_{n_k} = s, V_k] - (\mathcal{T}^\pi V_k)(s). \label{eq:noise_step1}
\end{align}
By Theorem~\ref{thm:mdp_expansion} (MDP expectation expansion) applied to the function $f(a, s') = r(s,a,s') + \gamma V_k(s')$:
\begin{align}
    \E[R_{n_k+1} + \gamma V_k(S_{n_k+1}) \mid S_{n_k} = s, V_k]
    &= \sum_a \pi(a|s) \sum_{s'} P(s'|s,a)\bigl[r(s,a,s') + \gamma V_k(s')\bigr] \nonumber\\
    &= (\mathcal{T}^\pi V_k)(s). \label{eq:noise_step2}
\end{align}
Substituting~\eqref{eq:noise_step2} into~\eqref{eq:noise_step1}: $\E[w_k(s) \mid \mathcal{F}_k] = (\mathcal{T}^\pi V_k)(s) - (\mathcal{T}^\pi V_k)(s) = 0$.

\textbf{Part 2.} Since the MDP is finite, rewards are bounded: $|R_{n_k+1}| \leq R_{\max}$. The value estimate satisfies $|V_k(s')| \leq \|V_k\|_\infty$. Therefore:
\begin{align}
    |w_k(s)| &\leq |R_{n_k+1}| + \gamma |V_k(S_{n_k+1})| + |(\mathcal{T}^\pi V_k)(s)| \nonumber\\
    &\leq R_{\max} + \gamma \|V_k\|_\infty + (R_{\max} + \gamma \|V_k\|_\infty) \nonumber\\
    &= 2R_{\max} + 2\gamma \|V_k\|_\infty. \label{eq:noise_bound_step}
\end{align}
Therefore $w_k(s)^2 \leq (2R_{\max} + 2\gamma \|V_k\|_\infty)^2 \leq C(1 + \|V_k\|_\infty^2)$ for $C = 8\max(R_{\max}^2, \gamma^2)$.
\end{proof}

\begin{remark}
Theorem~\ref{thm:noise_decomp} shows that the TD update fits precisely into the Robbins-Monro framework (Theorem 2.6 in the main article): $V_{k+1}(s) = V_k(s) + \alpha_k(s)[h(V_k)(s) + w_k(s)]$ where $h(V)(s) = (\mathcal{T}^\pi V)(s) - V(s)$ has unique root $V^\pi$, and $w_k(s)$ is zero-mean noise with bounded variance. The convergence then follows from stochastic approximation theory.
\end{remark}

%% ============================================================
%% SECTION 11: SUMMARY TABLE
%% ============================================================
\section{Quick Reference}
\label{sec:reference}

\begin{center}
\renewcommand{\arraystretch}{1.5}
\begin{tabular}{p{3.2cm} p{5.5cm} p{5.5cm}}
\toprule
\textbf{Result} & \textbf{Statement} & \textbf{Where in TD Article} \\
\midrule
Linearity & $\E[aX + bY] = a\E[X] + b\E[Y]$ & Splitting $G_t = R_{t+1} + \gamma G_{t+1}$ \\
Pull-out & $\E[h(Z)X \mid Z] = h(Z)\E[X \mid Z]$ & Extracting $V(s)$ from $\E[\delta_t \mid S_t]$ \\
MDP expansion & $\E_\pi[f(A,S') \mid S\!=\!s] = \sum_a \pi \sum_{s'} P \cdot f$ & Every Bellman derivation \\
Tower property & $\E[\E[X \mid Y]] = \E[X]$ & Bellman recursion, MDS $\Rightarrow$ zero mean \\
Total variance & $\Var(X) = \E[\Var(X|Y)] + \Var(\E[X|Y])$ & Bias-variance tradeoff \\
MDS condition & $\E[w_n \mid \mathcal{F}_{n-1}] = 0$ & All convergence proofs \\
Uncorrelated MDS & $\E[w_m w_n] = 0$ for $m \neq n$ & Noise variance bounds \\
Jensen's & $\varphi(\E[X]) \leq \E[\varphi(X)]$ & $\|\mathbf{P}_\pi \mathbf{v}\|_{d^\pi}$ contraction \\
\bottomrule
\end{tabular}
\end{center}

\end{document}
